% Generated draft: requires mathematical and source review.
\documentclass[11pt]{article}
\usepackage[margin=25mm]{geometry}
\usepackage{fontspec}
\setmainfont{Latin Modern Roman}
\usepackage{amsmath,amssymb,mathtools,unicode-math}
\setmathfont{Latin Modern Math}
\usepackage{enumitem,needspace,xurl,hyperref}
\setlength{\parindent}{0pt}
\setlength{\parskip}{5pt}
\setlength{\emergencystretch}{3em}
\newcommand{\N}{\mathbb{N}}
\newcommand{\Z}{\mathbb{Z}}
\newcommand{\Q}{\mathbb{Q}}
\newcommand{\R}{\mathbb{R}}
\newcommand{\C}{\mathbb{C}}
\newcommand{\F}{\mathbb{F}}
\newcommand{\A}{\mathbb{A}}
\newcommand{\PP}{\mathbb{P}}
\newcommand{\E}{\mathbb{E}}
\newcommand{\eps}{\varepsilon}
\newcommand{\dd}{\,\mathrm{d}}
\begin{document}
\section*{UnsolvedMath ID: 11000186}
\textbf{Source problem code:} catalogue entry
\subsection{Definitions and mathematical statement}
Put \(X=\operatorname{Hom}(\pi,G)/G\), with conjugation quotient as in the source. In the usual surface-group setting, \(\pi=\pi_1(\Sigma)\). Let \(\mathscr B\) denote the source's admissible bilinear forms \(B\), with associated symplectic forms \(\omega_B\). The standard smooth-locus construction pairs surface-group cohomology classes by cup product followed by \(B\) and evaluation on \([\Sigma]\); the exact hypotheses on \(G,B\) and the smooth stratum are absent.

The conjecture says that the \(\operatorname{Mod}_\Sigma\)-invariant subalgebra of measurable differential forms on \(X\) is generated by these \(\omega_B\):
\[
 \Omega^*_{\mathrm{meas}}(X)^{\operatorname{Mod}_\Sigma}
   =\text{the unital graded wedge subalgebra generated by }
          \{\omega_B:B\in\mathscr B\}.
\]
Thus invariants in all degrees, including invariant measurable functions in degree zero, are part of the statement. Whether one fixes a connected component or allows componentwise coefficients matters and is not specified. Measurable forms alone do not supply an ordinary de Rham differential without further regularity or weak-derivative conventions; the source's phrase de Rham algebra needs clarification. Equality and invariance pointwise versus almost everywhere are also unstated.

The continuation uses a finite measure \(\mu\) on \(X\) and \(\mathcal H=L^2(X,\mu)\), describes the mapping-class representation as unitary, and reports Andersen's communication that \(H^1(\operatorname{Mod}_\Sigma;\mathcal H)=0\). Unitarity of the usual pullback representation requires \(\mu\)-invariance, not finiteness alone; the missing measure definition must provide it. A further conjecture generalizing this one is announced but not included, and no higher-cohomology assertion is reconstructed.
\subsection{Short English statement}
The source conjectures that invariant measurable differential forms on the representation quotient are generated by its bilinear-form symplectic structures, with unresolved regularity, component and measure conventions.
\subsection{Sources}
\begin{itemize}
\item ``Mapping class group dynamics on surface group representations,'' in \emph{Problems on Mapping Class Groups and Related Topics}, edited by Farb, Conjecture 2.1 and continuation, supplied PDF page 215. \url{https://www.math.uchicago.edu/~farb/papers/mcgbook.pdf}. Recorded access 2026-07-29; extraction pdf-font-structure. The preceding definitions of \(G,\pi,B,\mu\) are absent.
\item Andersen, source-reported communication claiming the first-cohomology vanishing; no publication reference or complete attribution is supplied. The announced following conjecture is missing.
\item Local UnsolvedMath registry, \texttt{data/unsolvedmath/problems.json}, permanent ID \(11000186\), exact source code \texttt{catalogue entry}; supplied OPEN-TRIAGE summary says truncated.
\end{itemize}
\end{document}
